% ============================================================================
\section{Gradient Sparsity and Teacher Overlap}
\label{sec:subnetworks}
% ============================================================================

We ask whether OPD gradients concentrate on a small fraction of parameters,
whether different teachers affect the same parameters within one MOPD
student, and whether teacher distribution differences predict this overlap.
Each question includes its planned experiment; no outcome is assumed.

\subsection{Does gradient sparsity persist under multiple teachers?}

Prior work reports concentrated parameter changes in RL fine-tuning
\citep{mukherjee2025subnetworks} and OPD
\citep{yu2026denseupdates,shen2026opdgeometry}. These observations motivate
measuring both gradients and parameter changes in our setting.
For a domain gradient $g_d=\nabla_\theta\mathcal L_d$, where
$\mathcal L_d$ is that domain's mean loss, we report the fraction of parameters with
absolute gradient at most a threshold $\tau$, together with the fraction
of squared gradient norm carried by the largest-magnitude $\rho$ fraction
of parameters. We sweep $\tau$ and $\rho$: ``gradient sparsity'' here means
near-zero gradients or concentrated gradient energy, not necessarily
exact zeros.

\paragraph{Experiment.}
Using the representative single-teacher and joint runs from
Section~\ref{sec:density}, collect gradients at initialization and
early, intermediate, and final checkpoints. Compare each domain's gradient
and the combined MOPD gradient under the same loss reduction, reporting
curves against both domain exposure and compute. Match diagnostic response
counts per domain across single-teacher and joint students, report valid-token
counts, and check sensitivity to batch size; record the combined gradient's
full batch size separately. Repeat over fresh student-rollout batches,
record accumulation denominators, and
accumulate raw gradients in FP32 before clipping or optimizer processing.
Threshold curves, gradient norms, and energy concentration jointly show
whether apparent sparsity reflects small overall scale or concentration
on particular parameters.

As secondary measurements, record optimizer steps and cumulative changes
$\Delta=\theta-\theta_0$. These are not equivalent to gradients: AdamW
momentum and weight decay can move parameters with a zero current gradient,
and changes from successive steps can cancel. Compare FP32 master-weight
differences with exported BF16 checkpoint differences to identify
rounding-induced zeros; converting BF16 values to FP32 cannot recover
lost changes. This comparison checks the reported parameter-change
phenomenon while separately testing gradient sparsity in MOPD.

\subsection{Do teachers affect the same parameters in one student?}

At a fixed MOPD checkpoint $\theta$, compute each routed domain's gradient
$g_d$ using its teacher $\pi_{\phi_d}$ and student-generated responses.
We call parameters with large absolute gradients \emph{active parameters}
for that teacher and batch. This is a measurement of local sensitivity,
not a claim that the model contains a fixed, isolated subnetwork.
To compare locations fairly, first select the same fraction of parameters
for every teacher, ranked by absolute gradient. For domains $d$ and $e$,
\begin{equation}
\text{overlap}(d,e)=
\frac{\text{number of parameters selected by both teachers}}
     {\text{number of parameters selected by either teacher}}.
\label{subnet_overlap}
\end{equation}
Zero means no shared selected parameters; one means identical selections.
The corresponding difference is one minus this overlap. We also report
absolute-threshold selections, which retain information about gradient
scale as well as location.

\paragraph{Experiment.}
Every teacher calculation starts from the same frozen student checkpoint.
Repeat over diagnostic batches and report overlap matrices across
checkpoints, with concentration curves alongside them. Compare global
selections and selections made separately within each layer against
random selections with matching counts, including per-layer results.
Signed gradient alignment provides additional context: shared parameters
can receive compatible or opposing gradients, so lower overlap is not
assumed better.

Separately compare proposed optimizer steps, each computed from an
identical copy of the optimizer state without advancing training. Their
overlap can reflect shared momentum or decay and does not decompose the
joint AdamW step into additive teacher contributions. Independently trained
single-teacher endpoint changes provide a reference only: those students
follow different trajectories and cannot identify teacher contributions
inside the joint student. Appendix~\ref{app:records} specifies recording
and threshold conventions.

\subsection{Do more different teachers affect more different parameters?}

We measure the difference between teachers $\pi_{\phi_d}$ and
$\pi_{\phi_e}$ by their Jensen--Shannon divergence, a symmetric, bounded
measure of how differently they distribute probability over next tokens.
Every teacher scores the same cached, held-out student prefixes, with
common domain weights. We also record each teacher's difference from the
current student $\pi_\theta$ and its gradient and step norms.
Appendix~\ref{app:teacher_distance} gives the distribution-distance formula.

\paragraph{Experiment.}
At each checkpoint, plot teacher distribution difference against the
parameter-selection difference from the routed batches above. A smaller
control recomputes all teacher gradients on identical prefixes, testing
whether an association persists when the response contexts also match.
For PG, pair the same freshly sampled next tokens and fixed prefix weights
across teachers, repeating the paired draws across batches.
Available intermediate teacher checkpoints can add variation without
training new teachers. Code and science share one Qwen3-4B teacher, so
their routed comparison reflects context differences rather than two
independent teacher models.

Show individual pair trajectories and run-level uncertainty. Pairs that
share a teacher or training lineage, and repeated checkpoints, are
dependent observations. Relate teacher--student differences to gradient
concentration and observed transfer as additional analyses. A small
teacher pool supports a setting-specific association, not a universal
monotone relationship; larger, smaller, or unchanged parameter overlap
are all valid outcomes.
